\clearpage
\fancyhead[L]{\small GAUSS TO BODE.}
\hypertarget{native-419-part-lower}{}\label{native-419-part-lower}
\begin{center}
\Large GAUSS TO BODE.
\end{center}

\begin{center}
\rule{6cm}{0.4pt}

\textit{Astronomisches Jahrbuch für 1821}. pp. 212..216. Berlin, 1818.
\end{center}

\begin{flushright}Göttingen, 7 September 1818.\end{flushright}

\par\indent At your request, my much esteemed friend, I enclose another ephemeris for Pallas, from August 1819 to June 1820. It was calculated by Mr. \textsc{Dirksen}, who devotes himself zealously to the mathematical sciences here. He has also repeated the calculation of this year’s opposition using the same elements that Dr. \textsc{Tittel} used, and has found a somewhat different result:

\begin{center}
\setlength{\tabcolsep}{3pt}
\begin{tabular}{l|c|c}
 & Mr. Dirksen & Dr. Tittel \\
\hline
Time in Göttingen . & \shortstack{Sept. 8\\$15^{\mathrm u}\;17^{\mathrm m}\;5^{\mathrm s}$} & \shortstack{Sept. 8\\$18^{\mathrm u}\;44^{\mathrm m}\;12^{\mathrm s}$} \\
True longitude . . & $345^\circ\;47'\;18''$ & $345^\circ\;55'\;34''$ \\
Geocentric latitude & $6\;47\;25$ & $6\;47\;11$ \\
\end{tabular}
\end{center}

\par\indent I cannot say how this error in calculation arose, but it is confined to the calculation of the opposition, since several positions in the \SourceLeaf{420} ephemeris recalculated by Mr. \textsc{Dirksen} agree perfectly with Dr. \textsc{Tittel}’s figures (whose correctness is already confirmed by the observations I give below).

\par\indent I have been observing since the end of April this year with \textsc{Repsold}’s excellent meridian circle, an instrument you have doubtless seen for yourself before. It has now received so many important improvements, however, that it can be considered an entirely new instrument. It has been given a completely new graduation, three readings in place of the former single reading, and new pivot bearings of rock crystal. The pivots themselves have been turned again and centered, and so on. I am still expecting several other new fittings for it from Mr. \textsc{Repsold}. The following determinations of the right ascension of the North Star may serve as an example of the excellence of the instrument as a transit telescope. The last column, headed Weight, express the square of the precision according to the number of contributing threads, taking as the unit the precision of a determination made from one thread crossing at upper culmination and one at lower culmination. The formula for it is
\[
\frac{2ab}{a+b}
\]
\noindent if the determination is based on two culminations and $a$ and $b$ denote the numbers of crossings observed in each, and
\[
\frac{8abc}{ab+4ac+bc}
\]
\noindent for a determination from three successive culminations, with $c$ thread crossings observed in the third. In the first case one assumes that the piers have remained in place for 12 hours. In the latter one assumes only that any changes in them have been uniform over 24 hours. The times of each determination of right ascension are counted from the lower culmination of each day.

\begin{center}
\begin{tabular}{l|r|c|c}
1818 & \shortstack{Observed right asc.\\of the North Star} & \shortstack{Correction to\\\textsc{Bessel}’s table} & Weight \\
\hline
June 3.25 & $0^{\mathrm u}\,55^{\mathrm m}\,58^{\mathrm s}65$ & $+3^{\mathrm s}93$ & 2.67 \\
4.25 & 58.19 & $+2.80$ & 5.00 \\
5.25 & 58.56 & $+2.45$ & 5.00 \\
6.25 & 58.91 & $+2.22$ & 4.00 \\
8.50 & $0\,\;56\;\;1.14$ & $+2.97$ & 6.40 \\
10.25 & 2.58 & $+3.20$ & 3.75 \\
11.25 & 2.68 & $+2.59$ & 1.67 \\
12.50 & 3.31 & $+2.36$ & 4.29 \\
16.75 & 6.11 & $+2.13$ & 4.44 \\
18.25 & 7.31 & $+2.26$ & 5.00 \\
19.25 & 8.29 & $+2.51$ & 2.86 \\
21.25 & 10.49 & $+3.26$ & 5.00 \\
22.75 & 12.02 & $+3.60$ & 1.33 \\
26.25 & 13.49 & $+2.60$ & 1.60 \\
27.25 & 14.77 & $+3.14$ & 4.44 \\
29.50 & 16.71 & $+3.42$ & 3.43 \\
July 7.50 & 22.51 & $+3.31$ & 6.00 \\
8.25 & 22.97 & $+3.22$ & $1\;67$ \\
9.25 & 23.89 & $+3.40$ & 3.75 \\
Sept. 3.75 & 57.44 & $+2.35$ & 2.86 \\
\end{tabular}
\end{center}

\par\indent If one assumes that the correction to \textsc{Bessel}’s table is constant throughout this period, its mean would be
\[
= +2^{\mathrm s}86
\]

\Needspace{20\baselineskip}\SourceLeaf{421}
\par\indent I also add my observations of Uranus, Saturn, and Pallas. The oppositions of the last two planets are still a few days away.

\begin{center}
\emph{Observations of Uranus.}

\setlength{\tabcolsep}{2.2pt}\begin{tabular}{l|ccc|ccc|ccc}
\hline
1818 & \multicolumn{3}{c|}{Mean time} & \multicolumn{3}{c|}{Right asc.} & \multicolumn{3}{c}{Declination} \\
\hline
June 3 & $12^{\mathrm u}$ & $22^{\mathrm m}$ & $17^{\mathrm s}5$ & $257^\circ$ & $23'$ & $28''3$ & $23^\circ$ & $1'$ & $41''6$ South \\ 
4 & 12 & 18 & 10.9 & & 20 & 47.8 & & 1 & 29.0 \\
5 & 12 & 14 & $4\;6$ & & 18 & 9.9 & & 1 & 14.4 \\
7 & 12 & 5 & 51.6 & & 12 & 54.6 & & 0 & 50.5 \\ 
8 & 12 & 1 & 45.1 & & 10 & 15.0 & & 0 & 38.4 \\
9 & 11 & 57 & 38.6 & & 7 & 36.4 & & 0 & 27.4 \\
10 & 11 & 53 & $32\;1$ & & 4 & 56.5 & & 0 & 12.6 \\
11 & 11 & 49 & 25.7 & & 2 & 18.0 & & 0 & 2.9 \\
12 & 11 & 45 & 19.2 & 256 & 59 & 33.9 & 22 & 59 & 50.9 \\
13 & 11 & 41 & 12.7 & & 57 & 0.4 & & 59 & 40.5 \\
\hline
\end{tabular}
\end{center}

\begin{center}
\emph{Observations of Saturn.}

\setlength{\tabcolsep}{2.2pt}\begin{tabular}{l|ccc|ccc|ccc}
\hline
Aug. 31 & $12^{\mathrm u}$ & $31^{\mathrm m}$ & $25^{\mathrm s}0$ & $347^\circ$ & $24'$ & $5''2$ & $7^\circ$ & $48'$ & $6''2$ South \\
Sept. 1 & 12 & 27 & 12.0 & & 19 & 49.3 & & 49 & 55.8 \\
3 & 12 & 18 & 47.0 & & 11 & 30.1 & & 53 & 37.4 \\
4 & 12 & 14 & 34.1 & & 7 & 14.1 & & 55 & 27.4 \\
6 & 12 & 6 & 8.7 & 346 & 58 & 50.2 & & 59 & 7.0 \\
\hline
\end{tabular}
\end{center}

\begin{center}
\emph{Observations of Pallas.}

\setlength{\tabcolsep}{2.2pt}\begin{tabular}{l|ccc|ccc|ccc}
\hline
Aug. 31 & $12^{\mathrm u}$ & $25^{\mathrm m}$ & $8^{\mathrm s}7$ & $345^\circ$ & $49'$ & $45''1$ & $2^\circ$ & $23'$ & $2''4$ North \\
Sept. 4 & 12 & 6 & 26.8 & & 5 & 5.8 & 1 & 32 & 43.4 \\
\hline
\end{tabular}
\end{center}

\par\indent Observations of celestial bodies as faint as Pallas are made much more difficult by the fact that the extremely fine threads (the natural spider threads have been split several times), which are otherwise of such excellent service, especially for daylight observations, can scarcely be made visible at the level of illumination those bodies can tolerate. In future I shall fit separate, thicker threads for such purposes.

\Needspace{22\baselineskip}\par\indent Mr. \textsc{Dirksen} compared the obsorvations of Uranus with \textsc{de Lambre}’s tables of Uranus and found the following differences: 

\begin{center}
\begin{tabular}{l|c|c}
 & Right asc. & Declination \\
\hline
June 3 & $-54''1$ & $+10''3$ \\
4 & $-52.1$ & $+10.6$ \\
5 & $-52.9$ & $+12.9$ \\
7 & $-55.8$ & $+12.3$ \\
8 & $-55.9$ & $+12.0$ \\
\end{tabular}
\par\medskip
\begin{tabular}{l|c|c}
 & Right asc. & Declination \\
\hline
June 9 & $-56''8$ & $+10''5$ \\
10 & $-56.4$ & $+12.9$ \\
11 & $-57.2$ & $+10.1$ \\
12 & $-52.4$ & $+9.4$ \\
13 & $-58.1$ & $+7.4$ \\
Mean & $-55.2$ & $+10.8$ \\
\end{tabular}
\end{center}

\Needspace{16\baselineskip}\par\indent From this, Mr. \textsc{Dirksen} calculated the opposition as follows:

\begin{center}
\begin{tabular}{ll}
\multicolumn{2}{l}{9 June 1818. $5^{\mathrm u}\,30^{\mathrm m}\,43^{\mathrm s}$} \\
\multicolumn{2}{l}{Mean time in Göttingen.} \\
True longitude . . . . . . . . . . . . . & $258^\circ\;10'\;34''5$ \\
\end{tabular}\par\smallskip
\begin{tabular}{ll}
Geocentric latitude . . . . . . . . . . . & $0\;4\;12.1$ \\
Heliocentric . . . . . . . . . . . . . & $0\;3\;58.7$ \\
\end{tabular} $\left.\rule{0pt}{2.2em}\right\}$ South
\end{center}

\noindent and the error in those tables in heliocentric coordinates

\begin{center}
\begin{tabular}{ll}
Longitude . . . . . . . . . . . . . . . . & $-45''3$ \\
Latitude . . . . . . . . . . . . . . . . & $+14.6$ \\
\end{tabular}
\end{center}

\par\indent Only some of my numerous observations of fixed stars have been reduced so far, and I will report the results at a later date.
